\section{Related Work}

\subsection{Conflicting evidence and evidence-reducing RAG}

Retrieval-augmented generation conditions a language model on passages retrieved for the query \cite{lewis2020rag}. Because retrieved passages can be irrelevant, outdated or wrong, much work adds modules that decide which passages to trust; Xu et al.\ \cite{xu2024survey} survey conflicts between contexts and between context and the model's memory. Corrective RAG \cite{yan2024crag} grades each retrieved document and acts on the grade; Astute RAG \cite{wang2025astute} consolidates internal and external knowledge by source reliability; MADAM-RAG \cite{wang2025madam} has one agent per document debate, and introduces RAMDocs, which mixes ambiguity, misinformation and noise. ConflictRAG \cite{conflictrag2026} types each detected conflict and resolves factual ones by Entropy-TOPSIS credibility selection, temporal ones by recency and opinion conflicts by presenting several perspectives.

Two recent results are close in spirit to ours. CONFLICTS \cite{cattan2025conflicts} argues that the right behaviour depends on the conflict type, and specifies a different desired response for each. Ragability \cite{gross2026ragability} finds that warning a model of contradictions helps it detect them but hurts when there are none. Kumar \cite{kumar2026verification} shows that gating chunks one by one with an entailment check is worse than not filtering on multi-hop questions, since a chunk that does not entail the answer alone can still be necessary. We study a third case in which reducing evidence can hurt: the passage a module removes or never retrieves is the one that states \emph{when} the answer holds.

\subsection{Conditional answers}

ConditionalQA \cite{sun2022conditionalqa} annotates every answer with the conditions under which it holds, on UK government guidance; human performance on conditional answers with their conditions is much lower than on the rest of the benchmark. Chain of Condition \cite{lin2024chaincondition} prompts a language model to identify the conditions in a document, build their logical relations, and verify and resolve them against the scenario; evaluated mainly with gold passages, it improves how a reader handles the conditions it is given, not whether the pipeline gives them. Our study is complementary: it holds the reader fixed and measures where conditions are lost; our verification pass (Section~\ref{sec:diagnostics}) is a simple instance of Chain of Condition's identify-and-verify step. Defeasible scope has also been studied within single legal provisions: Chen et al.\ \cite{chen2026sgdt} parse statutes into span-grounded deontic trees and name the failure of dropping a defeater \emph{silent scope omission}. Their setting keeps a rule and its exceptions in one provision; ours asks what happens when chunking puts them in different passages.

\subsection{Chunking and retrieval granularity}

What a retriever indexes changes what it finds: Dense~X \cite{chen2024densex} shows that the retrieval unit (passage, sentence or proposition) affects both retrieval and downstream answers, and DAPR \cite{wang2024dapr} builds a benchmark of passages that need their document's context to be retrieved, ConditionalQA among them, on which adding the document title to each passage helps. Work on legal texts reports that fixed-size chunking separates basic rules from their exceptions \cite{prior2026chunking}, but does not measure how often or with what consequence. Closest to our setting, Wan et al.\ \cite{wan2026retrieverfree} retrieve from all ConditionalQA guides split into paragraphs and report retrieval hit rates and answer exact match; conditions are not evaluated.

\subsection{Multi-hop evidence, sufficiency and defeasible inference}

A condition chunked apart from its answer resembles a missing hop in multi-hop retrieval, since two passages are needed together, but the resemblance is partial. In multi-hop questions the second passage is needed to reach the answer at all, so losing it shows up as a wrong answer; a condition qualifies an answer that can be reached without it, so losing it shows up only as missing scope, which answer-centric metrics do not register. Per-chunk gating removes necessary bridge passages on multi-hop benchmarks \cite{kumar2026verification}; we measure the analogous removal for conditions, whose cost answer accuracy does not show. Joren et al.\ \cite{joren2025sufficient} classify whether a retrieved context suffices for a definitive answer, and find that models often answer rather than abstain when it does not; a context that has lost a condition can still support a definitive answer, so it is insufficient in a sense such a check does not test. In the terms of defeasible reasoning \cite{chen2026sgdt}, ConditionalQA's conditions are defeaters written into the source; we ask whether a pipeline delivers them and whether the generator states them.

\textbf{Our position.} Table~1 sets the closest lines of work against this study. We found no prior work that measures against human annotation how often an answer's conditions are lost in chunked open retrieval, attributes the loss to its stages, or evaluates conflict-handling operators where discarding is and is not correct; the novelty claim is limited to this combination, not to any one of its parts.

\begin{table}[!htbp]\centering
\textbf{Table 1.} The closest lines of work and what this study adds.\\\vspace{5pt}
{\fontsize{9}{11}\selectfont
\begin{tabular}{L{2.5cm}L{2.85cm}L{4.4cm}L{4.4cm}}
\toprule
\textbf{Line of work} & \textbf{Representative work} & \textbf{What it measures} & \textbf{What this study adds} \\
\midrule
Conflict handling and evidence reduction & CRAG \cite{yan2024crag}, Astute RAG \cite{wang2025astute}, ConflictRAG \cite{conflictrag2026}, MADAM-RAG \cite{wang2025madam}, CONFLICTS \cite{cattan2025conflicts}, Ragability \cite{gross2026ragability} & whether modules remove irrelevant, false or conflicting passages, or respond to the type of conflict & the same evidence-reducing steps where a removed passage can be a valid condition, with deletion rates measured against human annotation \\
Multi-hop evidence and sufficiency & Kumar \cite{kumar2026verification}, Joren et al.\ \cite{joren2025sufficient} & whether bridge passages survive filtering; whether a context suffices to answer & the loss of passages that qualify an answer reachable without them, which answer accuracy does not register \\
Conditional answers & ConditionalQA \cite{sun2022conditionalqa}, Chain of Condition \cite{lin2024chaincondition}, Wan et al.\ \cite{wan2026retrieverfree} & readers' answers and conditions given gold passages; open retrieval scored without conditions & conditions scored under chunked open retrieval, stage by stage, with the reader held fixed \\
Defeasible and scoped rules & Chen et al.\ \cite{chen2026sgdt} & rules and their defeaters parsed within one legal provision; silent omission of a defeater & defeaters that chunking separates from the rule they qualify \\
Chunking and retrieval granularity & Dense X \cite{chen2024densex}, DAPR \cite{wang2024dapr}, Prior et al.\ \cite{prior2026chunking} & how the retrieval unit and document context affect retrieval; rule--exception splitting, described qualitatively & how often conditions are separated and at what cost, with context-preserving retrieval designs compared (Section~\ref{sec:retrieval}) \\
\bottomrule
\end{tabular}
}
\end{table}
